\documentclass[11pt,a4paper]{article}
\usepackage[margin=2cm]{geometry}
\usepackage{mathptmx}
\usepackage{enumitem}
\usepackage{xcolor}
\usepackage[hidelinks]{hyperref}
\definecolor{acc}{HTML}{1F3864}
\setlength{\parindent}{0pt}\setlength{\parskip}{5pt}
\pagestyle{empty}
\begin{document}
{\Large\bfseries\color{acc} Project idea and insights}\\[4pt]
\textbf{Where Should the Work Go? HW/SW Partitioning and DMA Design for TinyML Inference on a RISC-V SystemC Platform}\\
Abdelrahman Youssef \textperiodcentered{} COE 407/661 System on Chip, KFUPM \textperiodcentered{} \url{https://github.com/AbdelRahmanAtef1911/soc-accel-project}

\section*{\color{acc}The idea}
A small chip with a CPU and an AI accelerator has to decide what runs in hardware, what stays in software, and how data moves between them. The whole chip was built in SystemC: a RISC-V CPU, a bus, memory, and an accelerator whose timing comes from STONNE. All 48 layers of the MLPerf Tiny models were run on it in five steps, each moving more work into hardware: software only, a simple accelerator, a better fabric, native convolution in hardware, and smarter data movement. At every step the time was split into its parts to see where it goes.

\section*{\color{acc}The insights}
\begin{enumerate}[leftmargin=1.5em, itemsep=4pt]
\item \textbf{Speeding up what isn't slow does nothing.} A faster accelerator changed the total by 0.1\%, because the CPU work was 96--99\% of the time.
\item \textbf{Where the work runs matters more than how fast the engine is.} Moving convolution into hardware gave 1{,}389$\times$ over software, 42$\times$ more than the better fabric alone.
\item \textbf{Every fix moves the bottleneck.} First it was CPU work, then data movement (76--92\% of the time), then sharing the bus.
\item \textbf{Once compute is cheap, moving data is the cost.} A wider DMA gave 1.5--1.7$\times$; a faster compute fabric would have given almost nothing.
\item \textbf{Overlapping transfers with compute only hides the compute.} Double buffering gave up to 1.28$\times$, and only 4--13\% where the time is almost all transfers.
\item \textbf{An interrupt frees the CPU; it doesn't make the job faster.} It is 2\% slower than polling on a private bus. It wins when the bus is shared, because a polling CPU steals bus time from the DMA (+9.7\% vs +2.2\% with cycle stealing).
\item \textbf{How the DMA shares the bus decides who waits.} Block mode keeps the accelerator fast but freezes the CPU (up to 26{,}624 cycles). Transparent mode keeps the CPU free but slows inference by up to 80\%. Cycle stealing with large bursts gave the best balance: 1024-byte bursts cut accelerator time by 35\% while keeping each CPU stall to one burst.
\end{enumerate}
\end{document}
